% ============================================================
% paper/sections/02_ontological_grounding_oc.tex
% Ontological Grounding: Ontology of Continua (OC Context) (finalized)
% ============================================================

\section{Ontological Grounding: Ontology of Continua Context}

\subsection{Necessity of an Ontological Grounding}

The Enterprise Continuum $K_{\mathrm{EA}}$ is not defined as an isolated
or self-sufficient model.
It is an explicit instantiation within the Ontology of Continua (OC),
and its formal meaning is inseparable from that context.
Absent an explicit grounding, the structure of $K_{\mathrm{EA}}$ is
liable to be misread as metaphorical, framework-based, or taxonomic.

This section provides the minimal ontological context required for
correct interpretation.
It does not restate, modify, or extend OC Core, and it introduces no new
axioms or primitives.
Its function is strictly delimiting: to specify how $K_{\mathrm{EA}}$
inherits its meaning and constraints from OC.

\subsection{Continuum as a Primitive Object}

Within OC, a continuum is treated as a primitive formal object.
It is not defined by material composition, organizational form, or
semantic interpretation, but by the admissibility of its structural
primitives.

At minimum, a continuum is characterized by:
\begin{itemize}
  \item an admissible state space $\Omega$,
  \item a boundary $\partial\Omega$ delimiting that space,
  \item one or more axes $A$ representing distinguishable degrees of
        freedom,
  \item structural constraints expressed through thresholds $\Theta$.
\end{itemize}

Existence is therefore conditional.
A continuum exists only while these primitives remain well-defined and
mutually compatible.

\subsection{Dynamics, Cycles, and Persistence}

OC distinguishes between instantaneous admissibility and persistence
over time.
A continuum may satisfy all static admissibility conditions at a given
instant while failing to persist as a diagnosable object.

Persistence requires the presence of recurrent cycles $C$ that close
over an internal time scale $\tau$.
These cycles are not optimization trajectories or control loops; they
are minimal structural closures that prevent collapse to the boundary of
$\Omega$.

In the absence of such cycles, a continuum may exist transiently but
cannot be treated as a stable object of diagnosis.

\subsection{Continuity and Structural Connectivity}

Continuity in OC is defined in structural rather than geometric terms.
The continuity measure $k \in [0,1]$ quantifies the degree of connectivity
among admissible states under the defined dynamics.

A value of $k = 0$ indicates a hard structural discontinuity: either the
state space is empty or admissible connectivity has collapsed.
In this regime, the continuum ceases to be a valid object of diagnosis,
independently of any external interpretation.

\subsection{Thresholds and Structural Tension}

Thresholds $\Theta$ define the admissibility limits of axes, potentials,
and cycles.
Structural tension arises when internal dynamics approach or exceed
these limits.

Within OC, thresholds are neither targets nor normative benchmarks.
They are formal constraints separating admissible from non-admissible
regimes.
Crossing certain thresholds results in phase transitions or termination
of the continuum, rather than graded degradation.

\subsection{Parent Continua and Compatibility}

No continuum exists in isolation.
Each continuum is embedded within one or more parent continua that
constrain its admissible state space and dynamics.

Compatibility with the parent continuum is a necessary condition for
existence.
Violation of parent compatibility does not merely reduce performance or
stability; it invalidates the definition of the child continuum itself.

\subsection{Implications for Enterprise Diagnosis}

Applying OC to enterprise diagnosis entails that:
\begin{itemize}
  \item enterprises are not assumed to exist unconditionally,
  \item diagnosability is conditional on structural admissibility,
  \item termination of diagnosis is a formal outcome, not an analytical
        failure.
\end{itemize}

$K_{\mathrm{EA}}$ inherits these implications directly.
The enterprise continuum is therefore subject to the same ontological
constraints as any OC-defined continuum.

This grounding is required to interpret subsequent sections of this
whitepaper, particularly those concerning phase logic, admissibility
thresholds, and STOP conditions.
